\SourceLeaf{453}%
\Needspace{6\baselineskip}
\[
\dfrac{3014{,}48021}{3014{,}5757}=\operatorname{num}(\log=-0{,}00001376)
\]
\noindent for the baseline length according to \textsc{Peters} and the dimensions of the Earth stated by \textsc{Gauss} in the ‘Breitenbestimmung’, as given above,
\par

\Needspace{6\baselineskip}
\[
\dfrac{3014{,}451}{3014{,}5757}=\operatorname{num}(\log=-0{,}00001797)
\]
\noindent for the baseline length according to \textsc{Peters} and \textsc{Andrae} and the dimensions of the Earth stated by \textsc{Gauss} in the ‘Breitenbestimmung’, as given above,
\par

\Needspace{6\baselineskip}
\[
\dfrac{3014{,}48021}{3014{,}5757}.\dfrac{57009{,}746}{57009{,}7584}=\operatorname{num}(\log=-0{,}00001386)
\]
\noindent for the baseline length according to \textsc{Peters} and \textsc{Walbeck}’s dimensions of the Earth,
\par

\Needspace{6\baselineskip}
\[
\dfrac{3014{,}451}{3014{,}5757}.\dfrac{57009{,}746}{57009{,}7584}=\operatorname{num}(\log=-0{,}00001806)
\]
\noindent for the baseline length according to \textsc{Peters} and \textsc{Andrae} and \textsc{Walbeck}’s dimensions of the Earth,
\par

\Needspace{6\baselineskip}
\[
\dfrac{3014{,}48021}{3014{,}5757}.\dfrac{57009{,}746}{57008{,}551}=\operatorname{num}(\log=-0{,}00000466)
\]
\noindent for the baseline length according to \textsc{Peters} and \textsc{Schmidt}’s determination IV of the dimensions of the Earth,
\par

\Needspace{6\baselineskip}
\[
\dfrac{3014{,}451}{3014{,}5757}.\dfrac{57009{,}746}{57008{,}551}=\operatorname{num}(\log=-0{,}00000887)
\]
\noindent for the baseline length according to \textsc{Peters} and \textsc{Andrae} and \textsc{Schmidt}’s determination IV of the dimensions of the Earth
\par

\Needspace{6\baselineskip}
\[
\dfrac{3014{,}48021}{3014{,}5757}.\dfrac{57009{,}746}{57013{,}109}=\operatorname{num}(\log=-0{,}00003938)
\]
\noindent for the baseline length according to \textsc{Peters} and \textsc{Bessel}’s dimensions of the Earth,
\par

\Needspace{6\baselineskip}
\[
\dfrac{3014{,}451}{3014{,}5757}.\dfrac{57009{,}746}{57013{,}109}=\operatorname{num}(\log=-0{,}00004359)
\]
\noindent for the baseline length according to \textsc{Peters} and \textsc{Andrae} and \textsc{Bessel}’s dimensions of the Earth.
\par

The dimensions of the Earth calculated by \textsc{Schmidt} and by \textsc{Bessel} presuppose the length that \textsc{Schumacher} gave for his Braak baseline. Recalculating the arc measurements they considered using this corrected baseline length would give somewhat different figures for the dimensions of the Earth. Those figures will have to be replaced before long by better determinations as several new arc measurements are expected to be completed soon.

The coordinates specially included here in print from the partial registers are either the same as in the general register or based on less accurate determinations. they can, however, help explain the following ‘Outlines’. Of the partial registers, only Nos. 1 to 11 are in \textsc{Gauss}’s estate. A new comparison gave me the following corrections:
\begin{center}\small\begin{tabular}{@{}p{.39\linewidth}rrp{.16\linewidth}@{}}
In the general register:&\(-26619{,}9\)&\(-12689{,}5\)&Lauenberge\\
In partial register (3):&\(-26619{,}9\)&\(+12689{,}5\)&Lauenberge
\end{tabular}\end{center}
\begin{center}\begin{tabular}{@{}rrp{.55\linewidth}@{}}
\(-233491{,}171\)&\(+181317{,}782\)&Norden, turret on the tall church. No. 12.\\
\(-249451{,}172\)&\(+162280{,}007\)&Langeoog, F. J. \textsc{Paul}’s house, east gable pole. No. 12.
\end{tabular}\end{center}

To identify the points listed in the coordinate register more easily, one may make good use of ‘\textsc{Papen}’s map of the Kingdom of Hanover’, which is based on this survey.

I added the headings \(+\) south and \(+\) west to make it easier for the calculator to avoid misreading the signs.
\begin{flushright}\textsc{Schering.}\end{flushright}
\begin{center}\rule{.23\linewidth}{.3pt}\end{center}
